\section{Free Will, Leverage, and Decision Windows}
\label{sec:windows}

\subsection{Volitional Energy}

\begin{definition}[Volitional energy]
Volitional energy $\volition(t) \geq 0$ is the limited resource an agent uses
to override current-state inertia. It is finite within a session, depleted by
costly transitions, and partially restored by rest, closure, and recovery.
\end{definition}

\begin{definition}[Decision window]
A decision window is a short interval in which the effective transition
barrier is locally reduced:
\[
t \in \decisionwindow \quad \Longrightarrow \quad
\barrier(t) \leq \barrier(t + \varepsilon)
\]
for nearby $\varepsilon$ of ordinary size.
\end{definition}

\begin{theorem}[Leverage principle]
The practical influence of free will is maximized when high volitional energy
coincides with a low activation barrier.
\end{theorem}

\begin{discussion}
This theorem is straightforward but strategically important. People often
imagine freedom as a constant force that should defeat any obstacle on command.
The leverage principle says otherwise: freedom is most valuable when applied at
the right moment, not when exhausted against the worst possible timing.
\end{discussion}

\subsection{Four Classes of Decision Windows}

\begin{enumerate}
  \item \textbf{Morning reset windows.} The previous behavioural frame has not
        fully reloaded, so early inputs are disproportionately influential.
  \item \textbf{Boundary windows.} The completion of one task dissolves part of
        the previous context and temporarily lowers switching cost.
  \item \textbf{Environmental windows.} Movement into a library, office,
        classroom, or gym changes the cue structure before the old state is
        re-established.
  \item \textbf{Commitment windows.} An upcoming meeting, partner, or deadline
        compresses indecision and substitutes some external force for internal
        effort.
\end{enumerate}

\subsection{Missed Windows}

\begin{proposition}[Decay of leverage]
If an agent fails to act during a decision window, then the expected effort
required for the same transition weakly increases as the old frame re-stabilizes.
\end{proposition}

\begin{examplebox}
\textbf{Post-class transition.}

After finishing a lecture, the student has a narrow interval in which the next
best move is to walk directly to the study desk. If the student instead opens a
messaging app while standing in the hallway, the old academic frame dissolves
before the new study frame is established. The same one-hour study block now
feels more expensive despite no change in formal obligation.
\end{examplebox}

\subsection{Implications For Management}

If the theory is used in a coaching, education, or self-management setting, the
goal is not to maximize pressure. The goal is to design transitions so that
important actions occur near windows of low barrier. This is a cheaper and more
reliable control method than relying on repeated heroic effort.

\subsection{A Practical Diagnostic}

\begin{methodbox}[title={Window Diagnostic}]
When an important action fails to start, ask four questions:
\begin{enumerate}[nosep]
  \item Was a decision window available?
  \item If yes, was it recognized?
  \item If recognized, was the next action concrete enough to execute?
  \item If concrete, what re-stabilized the old frame before action occurred?
\end{enumerate}
\end{methodbox}

\begin{summarybox}
The chapter turns free will from a vague moral resource into a timing-sensitive
resource. That shift is the book's core economic claim: better timing and lower
barriers generate more reliable action than larger doses of abstract willpower.
\end{summarybox}

